\documentclass[11pt,a4paper]{article}
\usepackage[utf8]{inputenc}
\usepackage[margin=1in]{geometry}
\usepackage{amsmath,amssymb,amsthm}
\usepackage{listings}
\usepackage{xcolor}
\usepackage{hyperref}

\newtheorem{theorem}{Theorem}
\newtheorem{lemma}[theorem]{Lemma}

\lstset{
  language=C++,
  basicstyle=\ttfamily\small,
  keywordstyle=\color{blue}\bfseries,
  commentstyle=\color{gray},
  stringstyle=\color{red},
  numbers=left,
  numberstyle=\tiny\color{gray},
  breaklines=true,
  frame=single,
  tabsize=4,
  showstringspaces=false
}

\title{IOI 2015 -- Scales}
\author{}
\date{}

\begin{document}
\maketitle

\section{Problem Summary}

Six coins have distinct weights forming a hidden permutation of $\{1, 2, \ldots, 6\}$. Available queries (each with 3 outcomes):
\begin{itemize}
  \item \texttt{getLightest(A,B,C)}, \texttt{getMedian(A,B,C)}, \texttt{getHeaviest(A,B,C)}: return the lightest/median/heaviest of three coins.
  \item \texttt{getNextLightest(A,B,C,D)}: among $\{A,B,C\}$, return the lightest coin heavier than $D$ (wrapping cyclically).
\end{itemize}
Determine the complete ordering using at most a bounded number of queries.

\section{Solution}

\begin{theorem}
Six queries suffice. Since each query has 3 outcomes, 6 queries distinguish $3^6 = 729$ cases, which is at least $6! = 720$ permutations.
\end{theorem}

\textbf{Approach}: Build a \emph{decision tree} offline during \texttt{init}. Starting from the set of all 720 permutations, at each node choose the query that minimises the size of the largest branch (best 3-way split). Recurse until each leaf identifies a unique permutation.

During \texttt{orderCoins}, walk the precomputed tree, making queries at each internal node, until reaching a leaf.

\section{C++ Implementation}

\begin{lstlisting}
#include <bits/stdc++.h>
using namespace std;

struct Query { int type, a, b, c, d; };

int simulate(const Query &q, int perm[]) {
    int w[3] = {perm[q.a], perm[q.b], perm[q.c]};
    int c[3] = {q.a, q.b, q.c};
    for (int i = 0; i < 3; i++)
        for (int j = i+1; j < 3; j++)
            if (w[i] > w[j]) { swap(w[i],w[j]); swap(c[i],c[j]); }
    if (q.type == 0) return c[0];
    if (q.type == 1) return c[1];
    if (q.type == 2) return c[2];
    int wd = perm[q.d];
    for (int i = 0; i < 3; i++) if (w[i] > wd) return c[i];
    return c[0];
}

vector<vector<int>> allPerms;
vector<Query> allQueries;

void generatePerms() {
    vector<int> p = {1,2,3,4,5,6};
    do { allPerms.push_back(p); } while (next_permutation(p.begin(), p.end()));
}

void generateQueries() {
    for (int a = 0; a < 6; a++)
    for (int b = a+1; b < 6; b++)
    for (int c = b+1; c < 6; c++) {
        allQueries.push_back({0,a,b,c,-1});
        allQueries.push_back({1,a,b,c,-1});
        allQueries.push_back({2,a,b,c,-1});
        for (int d = 0; d < 6; d++)
            if (d!=a && d!=b && d!=c)
                allQueries.push_back({3,a,b,c,d});
    }
}

struct TreeNode {
    Query query;
    int children[6];
    int permIdx;
    bool isLeaf;
};

vector<TreeNode> tree;

int buildTree(vector<int> &cands, int depth) {
    if (cands.size() == 1) {
        TreeNode nd; nd.isLeaf = true; nd.permIdx = cands[0];
        memset(nd.children, -1, sizeof(nd.children));
        tree.push_back(nd);
        return (int)tree.size() - 1;
    }
    if (depth >= 8) return -1;

    Query bestQ; int bestMax = INT_MAX;
    map<int, vector<int>> bestSplit;
    for (auto &q : allQueries) {
        map<int, vector<int>> split;
        for (int idx : cands) {
            int perm[6];
            for (int i = 0; i < 6; i++) perm[i] = allPerms[idx][i];
            split[simulate(q, perm)].push_back(idx);
        }
        int mx = 0;
        for (auto &[k,v] : split) mx = max(mx, (int)v.size());
        if (mx < bestMax) { bestMax = mx; bestQ = q; bestSplit = split; }
    }

    TreeNode nd; nd.isLeaf = false; nd.query = bestQ; nd.permIdx = -1;
    memset(nd.children, -1, sizeof(nd.children));
    tree.push_back(nd);
    int nodeIdx = (int)tree.size() - 1;
    for (auto &[coin, sub] : bestSplit)
        tree[nodeIdx].children[coin] = buildTree(sub, depth + 1);
    return nodeIdx;
}

int rootIdx;

void init(int T) {
    generatePerms(); generateQueries();
    tree.clear();
    vector<int> all(720); iota(all.begin(), all.end(), 0);
    rootIdx = buildTree(all, 0);
}

void orderCoins() {
    int cur = rootIdx;
    while (!tree[cur].isLeaf) {
        Query &q = tree[cur].query;
        int res;
        if (q.type==0) res = getLightest(q.a+1,q.b+1,q.c+1)-1;
        else if (q.type==1) res = getMedian(q.a+1,q.b+1,q.c+1)-1;
        else if (q.type==2) res = getHeaviest(q.a+1,q.b+1,q.c+1)-1;
        else res = getNextLightest(q.a+1,q.b+1,q.c+1,q.d+1)-1;
        cur = tree[cur].children[res];
    }
    vector<int> &perm = allPerms[tree[cur].permIdx];
    vector<pair<int,int>> wc(6);
    for (int i = 0; i < 6; i++) wc[i] = {perm[i], i+1};
    sort(wc.begin(), wc.end());
    int W[6];
    for (int i = 0; i < 6; i++) W[i] = wc[i].second;
    answer(W);
}
\end{lstlisting}

\section{Complexity Analysis}

\begin{itemize}
  \item \textbf{Queries per test case}: at most 6 (information-theoretically tight, since $3^6 = 729 \ge 720 = 6!$).
  \item \textbf{Initialisation}: $O(720 \times |\text{queries}| \times 6)$, which is constant for fixed problem size.
  \item \textbf{Space}: $O(720)$ for the decision tree.
\end{itemize}

\end{document}
